% !PRINTED 229
\weberpage{229}
Außerdem setzen wir noch:
\weq{4}{P_{11}\,\eta(\omega)^{\frac{n-1}{2}} = e^{\frac{\pi i}{6}\mu'(\mu+\mu'\omega)\frac{n^2-1}{n}}\;\mathop{\Pi}\limits_{1,\frac{n-1}{2}}^{h}\vartheta_{11}(2h\varpi).}

Diese Größen $P$ lassen sich durch die Wurzeln der Teilungsgleichung ausdrücken, wenn man unter Anwendung von (3) die Quotienten
\[
\frac{P_{00}^2}{P_{01}P_{10}},\quad \frac{P_{10}^2}{P_{01}P_{00}},\quad \frac{P_{01}^2}{P_{10}P_{00}},\quad \frac{P_{11}^3}{P_{00}P_{01}P_{10}}
\]
bildet und vermittelst der Formeln des § 42 elliptische Funktionen einführt.

Man erhält so
\weq{5}{\begin{aligned}
(-1)^{\frac{n^2-1}{8}}\,2^{\frac{n-1}{2}}P_{00}^3 &= \mathop{\Pi}\limits_{1,\frac{n-1}{2}}^{h}\frac{\operatorname{dn}^2h\Omega}{\operatorname{cn}h\Omega},\\
(-1)^{\frac{n^2-1}{8}}\,2^{\frac{n-1}{2}}P_{10}^3 &= \mathop{\Pi}\limits_{1,\frac{n-1}{2}}^{h}\frac{\operatorname{cn}^2h\Omega}{\operatorname{dn}h\Omega},\\
(-1)^{\frac{n^2-1}{8}}\,2^{\frac{n-1}{2}}P_{01}^3 &= \mathop{\Pi}\limits_{1,\frac{n-1}{2}}^{h}\frac{1}{\operatorname{cn}h\Omega\,\operatorname{dn}h\Omega},
\end{aligned}}
\weq{6}{(-1)^{\frac{n^2-1}{8}}P_{11}^3 = (\varkappa\varkappa')^{\frac{n-1}{2}}\mathop{\Pi}\limits_{1,\frac{n-1}{2}}^{h}\frac{\operatorname{sn}^3h\Omega}{\operatorname{cn}h\Omega\,\operatorname{dn}h\Omega}.}

Diese Ausdrücke zeigen nun, daß $P_{00}^3$, $P_{10}^3$, $P_{01}^3$ \gesperrt{die Wurzeln von Transformationsgleichungen (Modulargleichungen) sind}. $P_{11}^6$ ist die Wurzel einer Multiplikatorgleichung, und wenn $n$ eine Quadratzahl ist, so ist auch $P_{11}^3$ \gesperrt{die Wurzel einer Multiplikatorgleichung}.

Man kann in mannigfaltiger Weise die Funktionen $P$ miteinander kombinieren, um neue Transformationsgleichungen zu erhalten. Wir führen folgende an:
\weq{7}{\begin{aligned}
\frac{P_{10}}{P_{00}} &= \mathop{\Pi}\limits_{1,\frac{n-1}{2}}^{h}\frac{\operatorname{cn}h\Omega}{\operatorname{dn}h\Omega},\\
\frac{P_{01}}{P_{00}} &= \mathop{\Pi}\limits_{1,\frac{n-1}{2}}^{h}\frac{1}{\operatorname{dn}h\Omega},
\end{aligned}}
%CONT>>
